%----------------------------------------------------------------------
\section{Structural Analogies}\label{app:analogies}
%----------------------------------------------------------------------

The following interpretive frameworks provide useful
intuition for the scheduler but are not load-bearing
for any theorem in the main text. They are collected
here for readers familiar with these areas.

\subsection{RISC ISA design}

The atom inventory of Section \ref{sec:atoms} is, in
form, a RISC instruction set. Each atom has fixed
arity, fixed latency (modulo batch-size scaling),
and a regular cost model. Compound productions play
the role of instruction fusion in the sense of
super-scalar CPU design: multiple atomic instructions
are dispatched together when their data dependencies
admit it. The analogy is close enough that
\texttt{fused\_matrix\_apply} corresponds roughly to
AVX-512 VPAND-VPXOR chains, and
\texttt{time\_share\_register} corresponds to
register-renaming across a pipeline stall boundary.

Where the analogy breaks down: the CSTZ planner
operates on $\GF(2)$-linear atoms exclusively for
fusion, whereas CPU super-scalar dispatch handles
arbitrary instructions by reordering under data
dependencies. Our ``linearity constraint'' corresponds
to the CPU's implicit reliance on architectural
registers being the only shared state: two
instructions can execute concurrently if and only if
their register live sets are disjoint.

\subsection{Software pipelining and kernel fusion}

Compiler-theoretic software pipelining (Muchnick 1997
Chapter 20) overlaps successive iterations of a loop
by staggering their scheduling. Kernel fusion in the
GPGPU literature (Sung et al.\ 2010) merges two GPU
kernels into one by interleaving their operations.
Our \texttt{time\_share\_register} is a degenerate
case of both: it pipelines $n$ distinct operations on
one register without loop structure, and it fuses
them into one kernel dispatch without crossing a
lane boundary.

\subsection{Pareto-A$^*$ search}

Mandow and P\'{e}rez-de-la-Cruz (2010) formulate
multi-objective shortest-path search as A$^*$ over a
partially ordered cost space. Their algorithm
generalizes uniform-cost A$^*$ by maintaining a
Pareto frontier of candidate paths per node rather
than a single best-cost-so-far. Our chart parser is
essentially their algorithm applied to the parse
forest: chart cells are nodes, Pareto-nondominated
items are candidate paths, and the admissible
heuristic is the component-wise minimum over
remaining workload atoms.

\subsection{Constrained parsing in machine translation}

The statistical machine translation literature
(Chiang 2005, Lopez 2008, Huang 2008) has developed
extensive Lagrangian-dual techniques for constrained
parse-forest search, motivated by decoding
constraints (language model integration, length
penalties, feature-weight optimization). Our
Section \ref{sec:lagrangian} imports this framework
with the substitution of the resource vector for
the feature vector. The subgradient iterations are
identical; only the interpretation of the dual
prices differs.

\subsection{Sheaves on sites}

[DAF]'s sheafification over Grothendieck topologies
and our chart parser are, at an abstract level,
instances of the same construction: starting from
local data (presheaf / workload atoms) and a
topology (site / production rules), extend to a
global section (sheaf / plan) by applying the
gluing axiom (intersection consistency /
cost-bounded completer). The analogy is developed
further in Section \ref{sec:discussion}, and
formalizing it as a functor between the two
categories is a natural direction for follow-up.

\subsection{Reed--Muller codes and the Plotkin construction}

Reed--Muller codes $\mathrm{RM}(r, m)$ are
recursively defined via the Plotkin construction
$\mathrm{RM}(r, m+1) = \mathrm{RM}(r, m) |
\mathrm{RM}(r-1, m)$, where $|$ is the
concatenation operation. Our encoding lattice
(Section \ref{sec:encoding}) exploits this
recursion: the zeta transform butterfly of
Definition \ref{def:zeta} is precisely the
Plotkin concatenation read as a change of basis.
The grade filtration of $\bigwedge V$ corresponds
to the order filtration of RM: a grade-$\leq r$
element is encoded losslessly into
$\mathrm{RM}(r, n)$'s information dimension.

\subsection{SPPF and the plan portfolio}

[DAF]'s SPPF retains multiple parse derivations as
packed alternatives; our portfolio retains multiple
plan derivations. The structural parallel runs
deep: in both cases, ambiguity is information, not
noise; in both cases, the data structure is
chosen to admit compact representation of
exponential ambiguity without enumerating. A
portfolio with $|\mathcal{P}| = k$ is the
scheduler's analogue of an SPPF packed node with
$k$ alternatives.
